\documentclass[11pt]{article}
\usepackage{icmlko}
\begin{document}

\title{못 세우는 게 아니라 셀 것이 없었다 --- 그리고 하나만은 진짜로 못 세운다}
\author{월드모델 실험실}
\date{노트 576 · 577 · 2026년 8월 5일}
\maketitle

\begin{abstract}
하루 전 노트 $575$(논문 $445$)가 ``\textbf{집 밖 짝 둘이 아래 절반에서
무너진다}''를 냈다. 왜인지 물었다 --- 잡음인가 바닥인가 전이 실패인가.
\textbf{짝이 자기 자료로 낼 수 있는 천장}(그 짝 축만으로 $5$겹 교차검증)을
절반마다 재니 갈렸다.

\begin{center}
\begin{tabular}{lrrl}
\hline
아래 절반 & 천장 & 전이 & \\
\hline
KR 만화 & $\mathbf{+0.0199}$ & $+0.0688$ & 신호가 없다 \\
CN 만화 & $\mathbf{-0.0352}$ & $+0.0768$ & 신호가 없다 \\
비게임 앱 & $\mathbf{+0.3271}$ & $+0.0722$ & \textbf{전이가 $22\%$ 만} \\
\hline
\end{tabular}
\end{center}

\textbf{두 짝이 같은 이유로 무너지는 게 아니었다.} KR $\cdot$ CN 은
\emph{그 짝 자기 모형도} 못 세운다 --- 아래쪽 라벨이 바닥에 붙어 셀 것이
없다. \textbf{앱만 천장이 살아 있는데 전이가 못 간다.}
\end{abstract}

\begin{figure}[t]
\centering
\includegraphics[width=\textwidth]{figs/f1.pdf}
\caption{\textbf{가} --- 천장도 같이 무너진다, 앱만 예외.
\textbf{나} --- 크기 탓이 아니다. \textbf{다} --- $575$ 를 가른다.}
\end{figure}

\section{바닥이지 잡음이 아니다 (노트 576)}

먼저 ``집 밖 라벨이 더 시끄러워서''를 죽였다. 모바일 아래 절반 라벨에
$f \times \mathrm{SD}$ 가우스 잡음을 넣어 보면

\begin{center}
\begin{tabular}{lrrrrr}
\hline
$f$ & $0$ & $0.5$ & $1.0$ & $2.0$ & $3.0$ \\
\hline
모바일 아래 절반 & $+0.4418$ & $+0.4088$ & $+0.3381$ & $+0.2243$ &
$\mathbf{+0.1629}$ \\
\hline
\end{tabular}
\end{center}

목표는 $+0.0705$ 인데 \textbf{$3 \times \mathrm{SD}$ 에서도 안 닿는다.}
잡음으로 설명하려면 비현실적인 양이 필요하다 --- \textbf{기각.}

대신 라벨의 생김새를 보면 바로 나온다.

\begin{center}
\begin{tabular}{lrrr}
\hline
아래 절반 & 라벨 SD & IQR & 서로 다른 값 \\
\hline
모바일(학습) & $0.651$ & $1.100$ & $43.6\%$ \\
KR 만화 & $\mathbf{0.051}$ & $\mathbf{0.088}$ & $25.4\%$ \\
비게임 앱 & $0.606$ & $1.061$ & $\mathbf{9.8\%}$ \\
\hline
\end{tabular}
\end{center}

\textbf{KR 아래 절반은 사실상 한 값이다}(전체 SD 의 $1/12$). 앱의
$0\!\sim\!40\%$ 창은 서로 다른 값이 \textbf{$4.7\%$} 뿐 --- 스무 행 중 한
행만 고유값이다. \textbf{순위를 매길 것이 없다.}

\section{천장을 절반마다 재면 확정된다 (노트 577)}

``신호가 없다''는 읽기를 직접 잰다 --- 그 짝 \emph{자기} 축으로 $5$겹
교차검증한 천장(노트 $475$ 장치)을 절반마다 건다. \textbf{행이 반이라
내려가는 몫}은 같은 크기의 \emph{무작위} 절반으로 뺀다.

\begin{center}
\begin{tabular}{lrrr}
\hline
천장 & 전체 & 무작위 절반 & \textbf{아래 절반} \\
\hline
KR 만화 & $+0.7152$ & $+0.6984$ & $\mathbf{+0.0199}$ \\
비게임 앱 & $+0.6177$ & $+0.5953$ & $+0.3271$ \\
CN 만화 & $+0.3619$ & $+0.3479$ & $\mathbf{-0.0352}$ \\
\hline
\end{tabular}
\end{center}

\textbf{행을 반으로 줄인 것 자체는 천장을 거의 안 깎는다}($-0.02$).
\emph{아래} 절반이라는 것이 깎는다. KR 은 $+0.7152 \to +0.0199$ 로
\textbf{$36$ 분의 $1$}이 된다.

\begin{quote}
\textbf{KR 아래 절반에서는 전이가 천장을 넘는다} --- 전이 $+0.0688$ 대
천장 $+0.0199$. 그 짝 자기 자료로도 못 하는 것을 통합 모형이 (조금) 한다.
\textbf{이것은 전이 실패의 반대다.}
\end{quote}

\section{앱만 다르다}

\begin{center}
\begin{tabular}{lrr}
\hline
전이 $\div$ 천장 & 아래 절반 & 위 절반 \\
\hline
KR 만화 & (천장 $\approx 0$) & $0.78$ \\
비게임 앱 & $\mathbf{0.22}$ & $0.72$ \\
CN 만화 & (천장 $< 0$) & $\mathbf{0.94}$ \\
\hline
\end{tabular}
\end{center}

앱 아래 절반은 \textbf{천장이 $+0.3271$ 로 살아 있는데} 전이가 $+0.0722$
다. \textbf{여기가 이 판에서 지금 제일 크게 열려 있는 자리다} --- 잡을 것이
있는데 안 잡고 있는 유일한 칸이다.

왜 앱만인지는 노트 $576$ 의 표가 말해 준다 --- 앱 아래 절반은 \emph{동점은
많지만}(서로 다른 값 $9.8\%$) \emph{폭은 넓다}(SD $0.606$). \textbf{계단형
라벨}이다. 계단 사이는 구별할 수 있고, 그 짝 축은 구별하는데 통합 모형이
못 한다.

\section{셋째 짝에서 재현된다}

CN 만화 $352$ 행(노트 $440$)을 같은 자로 넣었다.

\begin{center}
\begin{tabular}{lrr}
\hline
CN 만화 & 천장 & 전이 \\
\hline
전체 & $+0.3619$ & $+0.3094$ \\
\textbf{위 절반} & $+0.5623$ & $\mathbf{+0.5285}$ \\
아래 절반 & $-0.0352$ & $+0.0768$ \\
\hline
\end{tabular}
\end{center}

\textbf{위 절반에서 통합 모형이 그 짝 전용 모형의 $94\%$ 를 낸다.}
노트 $575$ 의 ``위 절반은 가려낸다''가 셋째 짝에서 선다.

\begin{quote}
\textbf{위 절반에서 통합 모형은 그 짝 전용 모형의 $72\!\sim\!94\%$ 를
낸다}(KR $0.78$ $\cdot$ 앱 $0.72$ $\cdot$ CN $0.94$). \emph{이것}이
제품 층에 적을 문장이고, ``$\rho$ 가 얼마''보다 정직하다 --- 천장을
같이 적으면 얼마나 남았는지가 보인다.
\end{quote}

\section*{내 숫자를 깎는다 --- 두 번째}

논문 $445$ 가 ``앱 위 절반은 학습 도메인과 $-0.0895$ 차뿐''을 실었다.
그 값은 \textbf{모바일을 연도로, 앱을 날짜로} 통제해 견준 것이었다.
맞춰서 다시 냈다.

\begin{center}
\begin{tabular}{lrr}
\hline
앱 $-$ 모바일 (위 절반) & 연도 통제 & 날짜 통제 \\
\hline
모바일 & $+0.3919$ & $+0.3919$ \\
비게임 앱 & $+0.2106$ & $+0.2108$ \\
\textbf{차} & $\mathbf{-0.1813}$ & $\mathbf{-0.1811}$ \\
\hline
\end{tabular}
\end{center}

\textbf{통제 변수는 범인이 아니었다} --- 편상관은 통제의 단조 변환에
불변이고($\mathrm{yr} = \mathrm{ASOF} - 10^{t}/365.25$), 실제로 소수점
넷째 자리까지 같다. 범인은 \textbf{노트 $575$ 의 모바일 위 절반
$+0.3032$ 가 재현되지 않는 것}이다(씨앗 여섯으로 $+0.3919$).

\begin{quote}
차는 $-0.0895$ 가 아니라 $\mathbf{-0.1811}$ --- \textbf{두 배}다.
결론은 살아남지만(위 절반에서는 학습 도메인과 견줄 만하다) \textbf{한
자리 세게 적었다.} 논문 $444$ 에 이어 이 연쇄에서 두 번째다.
\end{quote}

\section*{조항}

\begin{quote}
\textbf{집 밖 성적을 낼 때 그 짝의 천장을 같이 낸다.} ``$+0.07$ 이다''는
못 하는 것처럼 들리지만, 천장이 $+0.02$ 면 \emph{할 것이 없었다}는
뜻이고 천장이 $+0.33$ 이면 \emph{못 한 것}이다. \textbf{같은 숫자가 반대
뜻이다.}
\end{quote}

\section*{남은 것}

\begin{itemize}
\item \textbf{앱 아래 절반의 $0.22$} --- 천장 $+0.3271$ 이 살아 있다.
그 천장을 내는 축이 무엇인지 보면(짝 축 다섯 중 어느 것) 실마리가 나온다.
\textbf{지금 이 판에서 제일 큰 열린 자리다.}
\item \textbf{CN 을 다섯 자에 넣을까} --- 셋째 집 밖 짝이 서면 거울(노트
$426$)이 단단해진다. $352$ 행이라 무게는 작다.
\item \textbf{모바일 위 절반 $+0.3032$ 가 어디서 나왔나} --- 재현이 안 된다.
행 집합인지 자르는 규칙인지 안 봤다.
\end{itemize}

\end{document}
